\documentclass[11pt,a4paper]{article}
\usepackage[T1]{fontenc}
\usepackage[utf8]{inputenc}
\usepackage{mathpazo}
\usepackage{microtype}
\usepackage[margin=1.1in]{geometry}
\usepackage{amsmath,amssymb,amsthm}
\usepackage{booktabs,tabularx,array}
\usepackage{enumitem}
\usepackage[most]{tcolorbox}
\usepackage{titlesec}
\usepackage[hidelinks]{hyperref}
\usepackage{setspace}
\setstretch{1.08}

\titleformat{\section}{\normalfont\Large\bfseries}{\thesection}{1em}{}
\titleformat{\subsection}{\normalfont\large\bfseries}{\thesubsection}{1em}{}

\theoremstyle{plain}
\newtheorem{proposition}{Proposition}
\theoremstyle{definition}
\newtheorem{definition}{Definition}
\theoremstyle{remark}
\newtheorem{remark}{Remark}

\newtcolorbox{caution}[1][]{colback=gray!6,colframe=gray!55,boxrule=0.5pt,arc=1pt,
  left=6pt,right=6pt,top=4pt,bottom=4pt,fonttitle=\bfseries\small,title={Caution#1}}
\newtcolorbox{statusbox}[1]{colback=white,colframe=black!70,boxrule=0.6pt,arc=0pt,
  left=8pt,right=8pt,top=6pt,bottom=6pt,fonttitle=\bfseries,title={#1}}

\newcommand{\tier}[1]{\textup{\textsc{[#1]}}}
\newcommand{\Eone}{\tier{E1}}
\newcommand{\Etwo}{\tier{E2}}
\newcommand{\Ethree}{\tier{E3}}
\newcommand{\Efour}{\tier{E4}}
\newcommand{\GSH}{\mathrm{GSH}}
\newcommand{\GSSG}{\mathrm{GSSG}}

\title{\textbf{Precursor Is Not Supply}\\[0.4em]
\large Glutathione, Sulfur, and the Narrative Closure of Supplement Chains}
\author{Flyxion\\ \small Independent Researcher}
\date{24 September 2026}

\begin{document}
\maketitle

\begin{abstract}
\noindent
A short explanatory passage about glutathione, of a kind now routinely produced in conversation with assistants and wellness sources alike, compresses several correct facts and several incorrect inferences into a single persuasive chain: glutathione is the body's ``master antioxidant''; its active thiol ``grabs free radicals and heavy metals and escorts them out''; garlic is a sulfur donor; and dietary sulfur therefore flows through cysteine into glutathione, so that an existing supplement habit is revealed to have been ``already pointing there.'' This essay separates the chain into its links and grades each one. The biochemistry of glutathione is among the best established in physiology. The dietary inferences built on it are mostly not. The central error is the substitution of a precursor relation for a supply relation: the fact that $A$ is made from $B$ does not imply that more $B$ yields more $A$, since synthesis is regulated, compartmented, and saturable, and since humans cannot build cysteine from arbitrary dietary sulfur at all. The essay then widens the analysis beyond precursors. Using undenatured type II collagen (UC-II) as a contrasting case, it shows that dietary compounds can relate to the body in at least six distinct ways (composition, precursor, supply, induction, signal, and tolerance), each requiring different evidence. A further case, glandular supplements made from thyroid, adrenal, and pituitary tissue, shows the same error operating on risk rather than benefit: what a tissue contains says little about what eating it delivers, because oral activity depends on molecule class. The essay argues that the most common error in supplement reasoning is to establish one relation and claim another. It closes by treating the passage's rhetorical ending as the most instructive part of it, because a chain that confirms what its reader already does is the chain least likely to be checked.
\end{abstract}

\begin{statusbox}{Statement of Epistemic Status}
\small
\begin{enumerate}[leftmargin=1.4em,itemsep=2pt]
\item This is an essay in the analysis of claims, prepared by an independent researcher without clinical or pharmacological training. It has not been peer reviewed.
\item Nothing here is medical advice. It does not recommend starting, stopping, or changing any supplement, food, or medication. N-acetylcysteine in particular is a drug with clinical uses, contraindications, and interactions, and decisions about it belong with a clinician.
\item Biochemical claims tagged \Eone\ reflect textbook physiology. Claims about human supplementation outcomes rest on a smaller and more contested literature and are tagged accordingly.
\item References were compiled during drafting and should be verified against the primary literature before citation.
\end{enumerate}
\end{statusbox}

\noindent Evidentiary tiers follow the convention of earlier work: \Eone\ established; \Etwo\ plausible, with adjacent support; \Ethree\ open; \Efour\ rhetorical or phenomenological, with no evidentiary chain.

\section{The Specimen}

The passage under examination runs to four sentences. It identifies glutathione as a tripeptide of glutamate, cysteine, and glycine; locates its activity in the thiol group of cysteine; describes that thiol as capturing free radicals and heavy metals and conveying them out of the body; and then turns to its reader, observing that garlic is a sulfur donor, that sulfur becomes cysteine and cysteine becomes glutathione, that N-acetylcysteine is the standard supplement precursor, and that the reader's existing habits had therefore been converging on glutathione all along.

The passage is not careless. Its first half is largely accurate. Its second half is fluent, well connected, and wrong in instructive ways. Each step follows grammatically from the one before, and the final sentence rewards the reader for a choice already made. This combination, a correct foundation supporting an incorrect extension, capped by a confirmation, is common enough in contemporary health explanation to deserve a name. Here it will be called \emph{narrative closure}: the point at which a chain of reasoning stops because it has arrived somewhere satisfying rather than because it has been verified.

\section{What the Passage Gets Right}

\subsection{Structure}

Glutathione is $\gamma$-L-glutamyl-L-cysteinylglycine. It is indeed composed of glutamate, cysteine, and glycine, and its characteristic chemistry is indeed that of the cysteine thiol. \Eone

One detail is omitted, and it matters. The bond between glutamate and cysteine is not an ordinary peptide bond. It joins the \emph{side-chain} ($\gamma$) carboxyl of glutamate to the amino group of cysteine, rather than the $\alpha$-carboxyl used in proteins. This unusual linkage protects glutathione from most ordinary peptidases and means that it is not made on ribosomes but assembled by dedicated enzymes. Calling glutathione simply ``a three-amino-acid peptide'' is correct in composition and misleading in kind: it is a specialized metabolite, not a fragment of protein.

\subsection{Synthesis}

Glutathione is synthesized in the cytosol in two ATP-dependent steps:
\begin{align*}
\text{glutamate} + \text{cysteine} &\xrightarrow{\ \text{GCL}\ } \gamma\text{-glutamylcysteine},\\
\gamma\text{-glutamylcysteine} + \text{glycine} &\xrightarrow{\ \text{GS}\ } \GSH,
\end{align*}
where GCL is glutamate--cysteine ligase and GS is glutathione synthetase. The first step is rate-limiting and is subject to feedback inhibition by glutathione itself. Cysteine availability is commonly the limiting substrate. \Eone

\subsection{Function}

Glutathione is present in most cells at millimolar concentrations, highest in the liver. Its roles include serving as the reducing cofactor for glutathione peroxidases, which reduce hydrogen peroxide and lipid hydroperoxides; conjugation of electrophilic compounds by glutathione S-transferases, the first step of the mercapturic acid pathway by which many xenobiotics are excreted; maintenance of protein thiols; and participation in redox signaling. \Eone

In the peroxidase reaction two molecules of reduced glutathione are oxidized to glutathione disulfide:
\[
2\,\GSH + \mathrm{H_2O_2} \xrightarrow{\ \text{GPx}\ } \GSSG + 2\,\mathrm{H_2O},
\]
and glutathione reductase regenerates the reduced form at the expense of NADPH:
\[
\GSSG + \mathrm{NADPH} + \mathrm{H^+} \xrightarrow{\ \text{GR}\ } 2\,\GSH + \mathrm{NADP^+}.
\]
The system is therefore a cycle, not a consumable. Its capacity depends as much on NADPH supply, and hence on the pentose phosphate pathway, as on the quantity of glutathione present.

\subsection{N-acetylcysteine}

N-acetylcysteine (NAC) is a cysteine prodrug and a genuine glutathione precursor. Its best-established use is as the antidote for acetaminophen (paracetamol) overdose, in which a toxic metabolite, NAPQI, depletes hepatic glutathione faster than it can be resynthesized. Supplying cysteine in that setting restores synthesis and prevents liver injury. \Eone

This is the strongest case in the entire passage, and it is instructive precisely because of how special its conditions are. NAC works in overdose because demand has driven glutathione far below its regulated set point and cysteine has become acutely limiting. Section~\ref{sec:supply} returns to why that condition does not generalize.

\section{The Name}

The passage offers a small etymological gloss, locating ``thi-one'' in the thiol group. The gloss is intuitive and not quite right. In systematic chemical nomenclature the suffix \emph{-thione} denotes a thiocarbonyl group ($\mathrm{C{=}S}$), which glutathione does not contain. The name is historical rather than systematic: it was coined in the early 1920s, when the compound was isolated and initially misidentified as a dipeptide, by combining \emph{glut}amic acid with \emph{thio}, for sulfur. The sulfur is real. The ``-one'' is an accident of coinage.

The point is minor, but it has the same shape as the passage's larger errors. A true fact (the activity is in the thiol) is joined to a plausible-sounding derivation (the name encodes it) that happens to be false. The derivation is convincing because it is tidy.

\section{``Master Antioxidant''}

\subsection{A title, not a mechanism}

``Master antioxidant'' has no technical meaning. It is a promotional epithet that has migrated from supplement marketing into general explanation. Glutathione is one component of an antioxidant network that also includes superoxide dismutases, catalase, peroxiredoxins, the thioredoxin system, and dietary antioxidants such as vitamins C and E, several of which regenerate one another. No single member governs the others, and the relative importance of each varies by tissue, compartment, and oxidant.

\subsection{Radicals are not simply grabbed}

The phrase ``grabs free radicals'' suggests stoichiometric scavenging: one molecule of glutathione meets one radical and removes it. Glutathione can react directly with some radicals, but much of its protective work is enzymatic and indirect, feeding peroxidases that remove peroxides before they generate radicals, and it is recycled rather than spent. A model in which more glutathione simply captures more radicals misdescribes the system in the way that a model of a bank in which more paper money means more wealth misdescribes an economy.

\subsection{Redox state is a ratio}

The reducing power of the glutathione couple is conventionally described by its half-cell potential,
\[
E_h \;=\; E^{\circ\prime} \;-\; \frac{RT}{2F}\,\ln\frac{[\GSH]^2}{[\GSSG]},
\]
where the squared term reflects that two reduced molecules form one disulfide. The potential depends on the \emph{ratio} of reduced to oxidized forms and, because of the square, on absolute concentration as well. Cells regulate this state within ranges associated with proliferation, differentiation, and apoptosis. Pushing it in either direction is not a neutral improvement.

\subsection{Oxidants are also signals}

Reactive oxygen species are not merely damage. Hydrogen peroxide in particular acts as a signaling molecule, and oxidative signals participate in immune function and in adaptation to exercise. The broader antioxidant supplement literature is sobering on this point. Large randomized trials of $\beta$-carotene in smokers found increased, not decreased, lung cancer incidence \cite{atbc1994,caret1996}, and at least one controlled study reported that vitamins C and E blunted some metabolic benefits of exercise training \cite{ristow2009}. None of these concerns glutathione directly. They establish only that ``more antioxidant'' is not a reliably safe direction. \Eone\ for the trial results; their bearing on glutathione precursors specifically is \Ethree.

\begin{caution}
The trials cited concern specific compounds, doses, and populations. They do not show that glutathione precursors are harmful. They show that the premise behind the epithet, that antioxidant capacity is a quantity to be maximized, has failed empirical tests before.
\end{caution}
\section{The Sulfur Chain}

\subsection{Where cysteine's sulfur comes from}

The passage's second half rests on the chain
\[
\text{garlic} \;\to\; \text{sulfur} \;\to\; \text{cysteine} \;\to\; \text{glutathione}.
\]
The middle arrow fails for humans in its general form. Humans cannot incorporate inorganic sulfur, or sulfur in arbitrary organic compounds, into cysteine. The sulfur in human cysteine comes from two sources only: cysteine eaten directly in dietary protein, and methionine, an essential amino acid whose sulfur is transferred to a serine backbone through the transsulfuration pathway. \Eone

``Sulfur donor'' is therefore not a category that maps onto cysteine supply. A food may be rich in sulfur and contribute essentially nothing to the cysteine pool, if its sulfur is not in the form of cysteine or methionine. For most people eating ordinary amounts of protein, dietary protein dwarfs any plausible contribution from a seasoning.

\subsection{What garlic actually contains}

Garlic's characteristic chemistry lies in organosulfur compounds: alliin and its product allicin, diallyl sulfides, and, in aged preparations, S-allylcysteine. These are sulfur-containing, and some are cysteine derivatives. None is a significant source of cysteine for glutathione synthesis in the quantities garlic is eaten.

There is a more defensible version of the garlic claim, and it points in a different direction. Some garlic organosulfur compounds appear, mostly in cell and animal studies, to induce the expression of antioxidant and detoxifying enzymes, including enzymes of glutathione synthesis, through the Nrf2 signaling pathway. If garlic influences glutathione at all, the more plausible route is signaling, not supply: it would act on the regulator rather than filling the tank. \Etwo\ as a mechanism in model systems; \Ethree\ as a meaningful effect of dietary garlic in humans.

\begin{remark}
The two versions of the claim make different predictions. A supply mechanism predicts that effect scales with sulfur content. A signaling mechanism predicts that effect depends on specific compounds, and may saturate or reverse at high doses, as inducible stress responses often do. The passage's version is the one with less support.
\end{remark}

\section{Precursor Is Not Supply}\label{sec:supply}

\subsection{The general principle}

The passage's deepest error is not about garlic. It is the silent substitution of one relation for another.

\begin{definition}[Precursor relation]
$B$ is a precursor of $A$ if $A$ is synthesized from $B$.
\end{definition}

\begin{definition}[Supply relation]
$B$ supplies $A$ if increasing the availability of $B$ increases the steady-state amount of $A$.
\end{definition}

\begin{proposition}
The precursor relation does not imply the supply relation. Supply holds only when $B$ is limiting for the synthesis of $A$, and when regulation of that synthesis does not compensate for the change.
\end{proposition}

\subsection{A minimal model}

Let the rate of glutathione synthesis through its rate-limiting enzyme be written schematically as
\[
v \;=\; V_{\max}\,\frac{[\mathrm{Cys}]}{K_m + [\mathrm{Cys}]}\cdot\frac{1}{1 + [\GSH]/K_i},
\]
where the first factor is ordinary saturable kinetics in cysteine and the second represents feedback inhibition by the product. Steady state requires that synthesis balance consumption and export, $v = k\,[\GSH]$.

Two consequences follow. First, if $[\mathrm{Cys}] \gg K_m$, the first factor is near one and further cysteine changes almost nothing. Second, even when cysteine is limiting, any rise in glutathione increases the inhibitory term, so the steady state moves less than the substrate does. The system is built to hold glutathione near a set point, not to track intake.

\begin{caution}
The model omits transport, compartmentation, the regulation of GCL expression, and export of glutathione from the liver to the plasma. It is used only to illustrate why a precursor need not act as a supply. It is not a quantitative description of human glutathione metabolism.
\end{caution}

\subsection{Supply and induction act on different terms}

The minimal model also clarifies the garlic question from Section~5. A supply intervention changes $[\mathrm{Cys}]$, the argument of the saturating factor, and is therefore subject to both saturation and feedback. An inducing intervention, one that increases expression of glutamate--cysteine ligase, changes $V_{\max}$ instead, and so raises the ceiling itself rather than pushing against it. Formally, near steady state,
\begin{gather*}
\frac{\partial v}{\partial [\mathrm{Cys}]} \;\to\; 0 \quad\text{as } [\mathrm{Cys}] \gg K_m,\\[4pt]
\text{whereas}\quad
\frac{\partial v}{\partial V_{\max}} \;=\; \frac{v}{V_{\max}} \;>\; 0 \quad\text{at every substrate level.}
\end{gather*}
This produces a small irony. The mechanism the passage proposed for garlic, supply of sulfur, is the one the kinetics disfavor even if the chemistry had allowed it. The mechanism the passage did not mention, induction, is the one that escapes the saturation argument entirely. Whether garlic induces glutathione synthesis in humans at dietary doses remains \Ethree, but if it does anything, it is more likely to act on the enzyme than on the substrate.

\begin{caution}
Escaping saturation does not make induction beneficial. Inducible stress responses are regulated for reasons, and chronically driving them is not evidently desirable. The point is only that induction and supply are different relations, with different kinetic signatures and different evidentiary requirements.
\end{caution}

\subsection{Why NAC works in overdose}

The acetaminophen case fits the model exactly. In overdose, consumption by NAPQI conjugation drives glutathione far below its set point, feedback inhibition is released, and cysteine becomes acutely limiting. Under those conditions the precursor relation temporarily \emph{does} become a supply relation, and supplying cysteine saves the liver. The case demonstrates the conditions under which precursors matter. It does not show that they matter when those conditions are absent.

This is the same structure as the error examined in earlier work on collagen supplements, where the fact that collagen contains particular amino acids was taken to imply that eating collagen supplies them to tissue as collagen. In both cases a relation of composition is read as a relation of delivery.

\section{Oral Glutathione}

Because glutathione is a peptide, it was long assumed to be largely broken down before absorption, making oral glutathione ineffective. Human trials have reported conflicting results. One controlled trial found no significant change in markers of glutathione status after supplementation \cite{allen2011}; another, longer trial reported increases in glutathione stores in blood and other compartments \cite{richie2015}. The question is open, and the existence of an increase in a measured compartment would not by itself establish any downstream health effect. \Ethree

\section{``Escorts Them Out''}

Glutathione does bind certain metals through its thiol, and glutathione conjugation contributes to the biliary excretion of some metals and many xenobiotics. \Eone\ for the physiology.

The phrase ``escorts them out,'' however, invites an inference the physiology does not support: that raising glutathione accelerates detoxification in an otherwise healthy person, or that supplementation constitutes a ``detox.'' Excretion depends on transporters, conjugating enzymes, liver and kidney function, and the chemical form of the metal. Heavy metal poisoning is a medical condition treated with specific chelating agents under supervision. \Ethree\ for supplementation as a detoxification strategy.

\section{A Relation of a Different Kind: Undenatured Collagen}\label{sec:ucii}

\subsection{Why the case belongs here}

The argument so far has turned on the difference between making and supplying. A useful test of any such distinction is a case that sits outside both. Undenatured type II collagen, sold commercially as UC-II and typically extracted at low temperature from chicken sternal cartilage, is such a case. It is neither a precursor nor a supply for anything, and its proponents do not claim that it is.

Hydrolyzed collagen, examined in earlier work, is marketed on a building-block logic that closely parallels the glutathione passage: the body contains collagen, collagen is made of these amino acids, therefore eating collagen supplies the body's collagen. UC-II is marketed on an entirely different logic. Its triple-helical structure is deliberately preserved, and its claimed mechanism is \emph{oral tolerance}: native type II collagen epitopes are presented to gut-associated lymphoid tissue, inducing regulatory T cells that are proposed to dampen immune responses against the type II collagen of joint cartilage. \Eone\ for oral tolerance as an immunological phenomenon; \Etwo\ for its induction by oral type II collagen in animal models; \Ethree\ for clinically meaningful effects in human joint disease.

\subsection{Dose as a discriminating variable}

The clearest evidence that UC-II belongs to a different category is its dose. It is typically taken at around 40\,mg per day. Adult dietary protein intake is commonly on the order of tens of grams per day, so the fraction contributed as substrate is
\[
\frac{m_{\text{UC-II}}}{m_{\text{protein}}} \;\approx\; \frac{4\times 10^{-2}\ \text{g}}{7\times 10^{1}\ \text{g}} \;\sim\; 10^{-3}\text{--}10^{-4}.
\]
No substrate or supply argument survives a ratio of that size. The product is either acting through a mechanism in which quantity is irrelevant beyond a threshold, as antigen presentation is, or it is not acting at all. The dose alone therefore excludes the building-block reading before digestion enters the discussion.

\begin{remark}
This is a general heuristic, not specific to collagen. When a compound is claimed to act at a dose that is nutritionally negligible, the claim is implicitly one of signaling, induction, or immunology, and must be evaluated as such. When a compound is claimed to act as a building block, its dose must be commensurate with the quantity of material being built. Either mismatch is diagnostic.
\end{remark}

\subsection{Mechanism and indication}

The most serious difficulty with UC-II is not its mechanism but the fit between mechanism and indication. Oral tolerance to type II collagen was developed and is best supported in collagen-induced arthritis, an animal model of rheumatoid arthritis, which is autoimmune. Oral type II collagen was trialed in human rheumatoid arthritis during the 1990s, where the mechanism should have been most applicable. An early small trial reported encouraging results \cite{trentham1993}; subsequent larger trials were, on the whole, disappointing \cite{barnett1998}. The direct human test of the mechanism in its best-matched disease therefore did not establish a reliable effect.

UC-II is nonetheless marketed principally for osteoarthritis, which is conventionally described as degenerative and mechanical rather than autoimmune. The mismatch is real, though it should not be overstated. Osteoarthritis involves synovial inflammation, and immune responses to cartilage matrix components exposed by degradation have been described. The plausibility of a tolerance mechanism in osteoarthritis is thus intermediate: weaker than in rheumatoid arthritis, stronger than zero. The human osteoarthritis trials that exist are small, short, and predominantly industry-funded \cite{lugo2016}. \Ethree

\begin{caution}
None of this constitutes a recommendation for or against UC-II. It is examined as a case in the logic of claims. Joint pain has many causes, and decisions about its treatment belong with a clinician.
\end{caution}

\subsection{What the case shows}

UC-II is the rare supplement whose claimed mechanism is coherent, specific, and matched to its form and dose. That makes it more interesting than hydrolyzed collagen and does not make it effective. It demonstrates that ``collagen supplement'' names a product category, not a mechanism, and that two products sharing a name can differ in every property that matters: structure, dose, target system, and the kind of evidence that would confirm them. It also shows the discipline of this essay working in the other direction. The same analysis that removes a false mechanism from the glutathione passage identifies a real, distinct mechanism in UC-II, and then declines to promote it beyond what the evidence supports.

\section{Tissue Is Not Dose: Glandular Supplements}\label{sec:glandular}

\subsection{A second specimen}

A second explanatory passage, of the same general kind as the first, addresses a supplement regimen combining organ meats (liver, heart, kidney) with glandular products made from thyroid, adrenal, and pituitary tissue. It divides the regimen sensibly into two halves: the organ meats are nutrition, while the glandulars are, in its phrase, ``endocrine intervention wearing a food costume.'' It then ranks the glandulars by risk, placing pituitary highest on the grounds that pituitary tissue is rich in growth hormone and ACTH.

The division is sound and the phrase is apt. The ranking, however, commits the error this essay has been tracing, in a new direction. The first specimen read composition as supply in order to promise a benefit. This one reads composition as supply in order to assign a risk. The logic is the same in both: a tissue contains a substance, therefore eating the tissue delivers the substance.

\subsection{Oral activity depends on molecule class}

Whether a hormone contained in eaten tissue reaches the circulation in active form depends principally on what kind of molecule it is.

\begin{description}[style=nextline,leftmargin=1.2em,itemsep=4pt]
\item[Iodothyronines.] Thyroxine (T4) and triiodothyronine (T3) are small amino-acid-derived molecules that are absorbed intact after oral intake. This is how prescription levothyroxine is given. Thyroid tissue eaten in any form therefore delivers active hormone. \Eone
\item[Steroids.] Cortisol and related corticosteroids are small lipophilic molecules that are orally active; prednisone and hydrocortisone are routinely given by mouth. Adrenal tissue containing corticosteroids delivers them. \Eone
\item[Catecholamines.] Adrenaline (epinephrine) and noradrenaline are small, but are extensively inactivated in the gut wall and liver, which is why adrenaline is given by injection or inhalation rather than swallowed. Adrenal medullary content therefore contributes little systemic catecholamine when eaten. \Eone
\item[Peptide and protein hormones.] Growth hormone, ACTH, TSH, prolactin, and the gonadotropins are proteins or peptides, and are degraded by digestion like other dietary protein. Growth hormone must be injected for this reason, as must insulin. Pituitary tissue eaten orally delivers these hormones as amino acids. \Eone
\end{description}

\subsection{A minimal expression}

The delivered systemic dose of a hormone $h$ from a glandular product may be written schematically as
\[
d_h \;=\; m \cdot c_h \cdot F_h, \qquad F_h \;=\; f_{\text{abs}}\, f_{\text{gut}}\, f_{\text{hep}},
\]
where $m$ is the mass of tissue consumed, $c_h$ the concentration of the hormone in that tissue after processing, and $F_h$ its oral bioavailability, decomposed into the fractions absorbed, escaping gut-wall metabolism, and escaping first-pass hepatic metabolism. The second specimen ranked products by $c_h$. Risk tracks $d_h$. For peptide hormones $F_h \approx 0$, so a high $c_h$ contributes nothing; for thyroid hormones $F_h$ is substantial, so even modest and variable $c_h$ matters.

On this expression the ranking inverts. Thyroid glandular is the product most likely to deliver an active hormone dose, and analyses of commercially available thyroid supplements have found measurable T3 and T4 in many of them \cite{kang2013}, with thyrotoxicosis from such products reported clinically. Adrenal glandular comes next, since corticosteroid content is orally active, and analyses of ``adrenal support'' products have reported the presence of thyroid hormone and steroid compounds \cite{akturk2018}. Pituitary glandular is, as a hormonal matter, the least active of the three.

\subsection{Where pituitary risk actually lies}

The pituitary gland sits within the skull, close to brain tissue. Its relevant risk class is not hormonal but infectious. Prion diseases are transmitted through central nervous system tissue, and the historical record contains a sobering precedent: growth hormone extracted from human cadaveric pituitaries and given by injection transmitted Creutzfeldt--Jakob disease to a number of recipients before recombinant growth hormone replaced it. The route (injection), species (human), and scale of that episode differ greatly from oral consumption of processed animal pituitary, and the precedent establishes only that pituitary tissue can carry prions, not that glandular products do. Regulators in several jurisdictions restrict bovine central nervous system tissues in food and supplements for this general reason. \Eone\ for the historical transmission; \Ethree\ for any risk from oral glandular products specifically.

\subsection{Uncontrolled dose as the defining property}

The single property that most distinguishes glandular products from medicines is not that they contain hormones, since many medicines are hormones. It is that $c_h$ is neither standardized nor reliably labeled. A prescribed thyroid hormone is taken at a known dose and monitored by blood test. A thyroid glandular delivers an unknown and batch-variable dose of the same molecule, typically without monitoring. The combination of orally active content and uncontrolled dose is what makes a product ``endocrine intervention wearing a food costume.'' Tissue origin alone does not.

\begin{caution}
This section is an analysis of claims, not medical advice. Anyone taking thyroid or adrenal glandular products regularly may wish to discuss them with a clinician; thyroid function in particular can be checked with a simple blood test (TSH), and symptoms such as palpitations, heat intolerance, tremor, insomnia, or unexplained weight loss are reasons to do so promptly. Stopping steroid-containing products after prolonged use can also carry risks, which is a further reason for clinical guidance.
\end{caution}

\subsection{The food half, and where supply does apply}

The organ meats in the same regimen are foods, and here the composition-to-supply inference largely holds. Vitamins and minerals in liver, heart, and kidney are absorbed in forms the body uses, at quantities commensurate with requirements. Organ meats are among the most nutrient-dense foods available. \Eone

One qualification completes the argument of Section~\ref{sec:supply}. Glutathione is protected from excess intake by saturation and feedback. Preformed vitamin A, abundant in liver, is not protected in the same way: it is stored in the liver and accumulates, and chronic excess intake is hepatotoxic and, in pregnancy, teratogenic. In the terms of the minimal model, retinol storage behaves more like a system without an effective feedback term. Supply is real here, and for that reason it needs a ceiling. Frequent liver consumption, particularly when combined with desiccated liver or other vitamin A sources, is the point at which ``nutrient-dense'' requires quantification.

\begin{remark}
The contrast is instructive in both directions. Where regulation is strong, as for glutathione, supply arguments fail because the system compensates. Where regulation is weak, as for stored vitamin A, supply arguments succeed, and that success is itself the hazard. In neither case does the direction of a precursor or composition relation settle the question.
\end{remark}

\section{A Typology of Relations}

The two cases together suggest that the relation between an ingested compound and a bodily outcome can take at least six forms. They are listed in order of increasing distance between the chemistry of the input and the chemistry of the output.

\begin{center}
\small
\begin{tabularx}{\textwidth}{@{}>{\raggedright\arraybackslash}p{0.15\textwidth}>{\raggedright\arraybackslash}X>{\raggedright\arraybackslash}p{0.2\textwidth}>{\raggedright\arraybackslash}X@{}}
\toprule
\textbf{Relation} & \textbf{Meaning} & \textbf{Example} & \textbf{Evidence required} \\
\midrule
Composition & The output contains the input's constituents & Glutathione contains cysteine & Chemical structure only \\
Precursor & The output is synthesized from the input & Cysteine to glutathione & Pathway biochemistry \\
Supply & More input yields more output & NAC in acetaminophen depletion & Limitation, dose--response, absence of compensation \\
Induction & The input raises synthetic capacity & Proposed garlic effect via Nrf2 & Enzyme expression changes at achievable doses \\
Signal & The input instructs target cells & Proposed Pro-Hyp effects on fibroblasts & Receptor, exposure, cellular response in tissue \\
Tolerance & The input reprograms an immune response & UC-II oral tolerance & Immune readouts and clinical outcome in matched disease \\
\bottomrule
\end{tabularx}
\end{center}

The typology makes the core error of supplement reasoning precise. Evidence establishing a relation high in the table is repeatedly offered for a relation lower down, or a relation of one kind is offered for a claim of another. Composition is presented as though it were precursor; precursor as though it were supply; and in hydrolyzed collagen, composition is presented as though it were supply directly, skipping every intervening row. The glutathione passage makes two such moves in succession. The typology does not say which relations hold for any given product. It says which evidence would be needed to show that they do.

\begin{remark}
The rows are not mutually exclusive. A single compound may stand in several relations at once, and NAC, for instance, has been studied both as a supply and as a compound with direct thiol chemistry of its own. The typology classifies claims, not molecules.
\end{remark}

The glandular case fits the table without strain. For thyroid and adrenal glandulars the supply row applies directly, in the unusual form of an exogenous hormone supplied as such. For pituitary glandulars only the composition row applies to their hormones: the tissue contains them, and eating it delivers amino acids.

\section{Narrative Closure}

\subsection{The final sentence}

The most instructive sentence in the passage is its last: the observation that the reader's existing choices had been converging on glutathione all along. It adds no information. Its function is to close the chain by aligning it with what the reader already does, so that accepting the chain requires no change and confirms a prior decision.

A chain that ends in confirmation is less likely to be checked than one that ends in a demand. There is no friction at its terminus, nothing to resist, no cost to accepting it. The reader's own habit becomes the evidence, and the chain's weakest link, the step from sulfur to cysteine, is carried across by momentum from the strong links before it and the pleasant conclusion after it.

\subsection{Anatomy of the pattern}

The pattern has four parts, each visible in the specimen:
\begin{enumerate}[itemsep=3pt]
\item \textbf{A correct anchor.} Real, checkable biochemistry establishes credibility.
\item \textbf{A category slide.} A word with a narrow meaning (cysteine as a sulfur-containing amino acid) is widened into a loose one (``sulfur''), which then admits new members (``sulfur donors'').
\item \textbf{A relation slide.} ``Is made from'' becomes ``is supplied by.''
\item \textbf{A personal terminus.} The chain ends at something the reader already does.
\end{enumerate}

Each part is individually innocuous. Accurate anchors are good practice; loose categories are how people speak; precursor and supply are easily confused; and relating information to the listener is ordinary courtesy. The combination is what produces confident error.

\subsection{Why the pattern persists}

The pattern persists because every step is locally plausible and the only global check, asking whether the conclusion would still be reached if the reader's habits were different, is never prompted. A useful discipline is to reverse the terminus: if the reader had instead been eating onions, broccoli, or eggs, would the passage have concluded that the pieces were pointing there too? Almost certainly, since all are commonly described as sulfur-rich. A conclusion that would have been reached from any starting point is not evidence about the starting point.

\section{Evidence Map}

\begin{center}
\small
\begin{tabularx}{\textwidth}{@{}>{\raggedright\arraybackslash}X l@{}}
\toprule
\textbf{Claim} & \textbf{Tier} \\
\midrule
Glutathione is a tripeptide of glutamate, cysteine, and glycine, with a $\gamma$-glutamyl bond & \Eone \\
Its characteristic activity depends on the cysteine thiol & \Eone \\
Synthesis is rate-limited by GCL and feedback-inhibited by glutathione & \Eone \\
Glutathione supports peroxidase activity and xenobiotic conjugation & \Eone \\
NAC restores glutathione and prevents liver injury in acetaminophen overdose & \Eone \\
Human cysteine sulfur comes from dietary cysteine or methionine & \Eone \\
Glutathione is the ``master'' antioxidant & \Efour \\
Garlic organosulfur compounds induce glutathione-related enzymes & \Etwo\ (models) \\
Dietary garlic meaningfully raises human glutathione & \Ethree \\
Garlic supplies sulfur that becomes cysteine for glutathione & Unsupported \\
Oral glutathione raises body glutathione stores & \Ethree \\
Raising glutathione accelerates metal or toxin excretion in healthy people & \Ethree \\
Precursor supplementation raises glutathione outside depletion states & \Ethree \\
\midrule
Oral tolerance is a real immunological phenomenon & \Eone \\
Oral type II collagen induces tolerance in arthritis models & \Etwo \\
Oral type II collagen reliably benefits rheumatoid arthritis & Not supported \\
UC-II acts as substrate for cartilage & Excluded by dose \\
UC-II meaningfully benefits osteoarthritis & \Ethree \\
\midrule
Eaten thyroid tissue delivers active thyroid hormone & \Eone \\
Some thyroid and adrenal supplements contain active hormones & \Etwo \\
Eaten pituitary tissue delivers active growth hormone or ACTH & Not supported \\
Pituitary tissue can transmit prions (injected human extract) & \Eone \\
Oral glandular products carry meaningful prion risk & \Ethree \\
Chronic excess preformed vitamin A is toxic & \Eone \\
\bottomrule
\end{tabularx}
\end{center}

\section{Conclusion}

Glutathione deserves its reputation among biochemists and does not need the one it has acquired elsewhere. It is a regulated, recycled, compartmented system for managing redox state and conjugating reactive compounds, and it is among the better-understood systems in physiology. What it is not is a tank filled by sulfur, drained by radicals, and topped up by seasoning.

The specimen passage is useful because it is mostly right. Its errors do not come from ignorance of glutathione. They come from two small substitutions, sulfur for cysteine and precursor for supply, and from an ending that made checking feel unnecessary. Undenatured collagen, set beside it, shows the same analysis producing a different verdict: a coherent mechanism, correctly matched to form and dose, whose evidence nevertheless remains thin. Glandular supplements show the error inverted, assigning risk by what a tissue contains rather than by what survives ingestion, so that the most hormonally inert product was ranked most dangerous and the most active was not. Across these cases lies a typology of six relations by which something eaten may bear on something built, and most confident claims in this area can be located by asking which relation was shown and which was asserted. The general lesson is modest: when an explanation of the body ends by confirming what one already does, the ending is the part to distrust first.

\begin{thebibliography}{9}
\small
\bibitem{allen2011} Allen J, Bradley RD. Effects of oral glutathione supplementation on systemic oxidative stress biomarkers in human volunteers. \emph{Journal of Alternative and Complementary Medicine} (2011). \emph{[Verify.]}
\bibitem{akturk2018} Akturk HK, et al. Over-the-counter ``adrenal support'' supplements contain thyroid hormone and steroid-like compounds. \emph{Endocrine Connections} (2018). \emph{[Verify.]}
\bibitem{atbc1994} The Alpha-Tocopherol, Beta Carotene Cancer Prevention Study Group. The effect of vitamin E and beta carotene on the incidence of lung cancer and other cancers in male smokers. \emph{New England Journal of Medicine} (1994).
\bibitem{barnett1998} Barnett ML, et al. Treatment of rheumatoid arthritis with oral type II collagen: results of a multicenter, double-blind, placebo-controlled trial. \emph{Arthritis and Rheumatism} (1998). \emph{[Verify.]}
\bibitem{caret1996} Omenn GS, et al. Effects of a combination of beta carotene and vitamin A on lung cancer and cardiovascular disease. \emph{New England Journal of Medicine} (1996).
\bibitem{kang2013} Kang GY, et al. Thyroid hormone content of commercially available dietary supplements. \emph{Thyroid} (2013). \emph{[Verify.]}
\bibitem{lugo2016} Lugo JP, et al. Efficacy and tolerability of an undenatured type II collagen supplement in modulating knee osteoarthritis symptoms: a multicenter randomized, double-blind, placebo-controlled study. \emph{Nutrition Journal} (2016). \emph{[Verify; industry-funded.]}
\bibitem{richie2015} Richie JP, et al. Randomized controlled trial of oral glutathione supplementation on body stores of glutathione. \emph{European Journal of Nutrition} (2015). \emph{[Verify.]}
\bibitem{ristow2009} Ristow M, et al. Antioxidants prevent health-promoting effects of physical exercise in humans. \emph{Proceedings of the National Academy of Sciences} (2009). \emph{[Verify.]}
\bibitem{trentham1993} Trentham DE, et al. Effects of oral administration of type II collagen on rheumatoid arthritis. \emph{Science} (1993). \emph{[Verify.]}
\end{thebibliography}

\end{document}
